% W4i42_Cosmology.tex
% DFD 17-Agent Research Programme — Iteration 42 (i42)
% Agents 10 (H0/DESI DR2), 11 (CMB R31), 14 (Wide Binaries/EFE)
% Theme: Hubble tension update, CMB third peak, Newtonian wide binaries
% Date: 2026-03-23

\documentclass[11pt,a4paper]{article}
\usepackage[utf8]{inputenc}
\usepackage{amsmath,amssymb,amsfonts,amsthm}
\usepackage{hyperref}
\usepackage{booktabs}
\usepackage{longtable}
\usepackage{geometry}
\usepackage{xcolor}
\geometry{margin=2.5cm}

\newtheorem{theorem}{Theorem}
\newtheorem{proposition}{Proposition}
\newtheorem{lemma}{Lemma}
\newtheorem{corollary}{Corollary}
\newtheorem{definition}{Definition}
\newtheorem{remark}{Remark}

\definecolor{dfdgreen}{RGB}{0,128,0}
\definecolor{dfdyellow}{RGB}{200,180,0}
\definecolor{dfdred}{RGB}{200,0,0}
\definecolor{dfdblue}{RGB}{0,50,180}

\newcommand{\GREEN}{\textcolor{dfdgreen}{\textbf{GREEN}}}
\newcommand{\YELLOW}{\textcolor{dfdyellow}{\textbf{YELLOW}}}
\newcommand{\RED}{\textcolor{dfdred}{\textbf{RED}}}
\newcommand{\NEW}{\textcolor{dfdblue}{\textbf{NEW}}}
\newcommand{\Tr}{\operatorname{Tr}}
\newcommand{\CP}{\mathbb{C}P}

\title{DFD i42: Cosmology --- $H_0$ Tension, CMB Third Peak, Wide Binaries\\[6pt]
\large Agents 10, 11, 14\\[6pt]
\normalsize Iteration 11/20 of Quantum Gravity Cycle (i32--i51)}
\author{DFD 17-Agent Research Programme}
\date{2026-03-23}

\begin{document}
\maketitle
\tableofcontents
\newpage

%=============================================================================
\part{Agent 10: Hubble Tension and DESI DR2 Update}
\label{part:agent10}
%=============================================================================

\section{Mandate}

Update the Hubble tension and dark energy status with the latest DESI DR2 analyses, including the debate about whether DESI truly reveals dynamical dark energy.

\section{$H_0$ Measurement Landscape (March 2026 Update)}

\begin{center}
\begin{longtable}{p{4cm}lll}
\toprule
\textbf{Method} & \textbf{$H_0$ (km/s/Mpc)} & \textbf{Year} & \textbf{Tension with Planck} \\
\midrule
\endhead
Planck 2018 CMB & $67.36 \pm 0.54$ & 2018 & --- \\
\midrule
ACT DR6 + BAO & $68.22 \pm 0.36$ & 2025 & $1.6\sigma$ \\
\midrule
SH0ES (Cepheids + JWST) & $73.0 \pm 1.0$ & 2025 & $5.0\sigma$ \\
\midrule
TRGB (Freedman/CCHP) & $69.8 \pm 1.7$ & 2024 & $1.2\sigma$ \\
\midrule
JWST Cepheids & $72.6 \pm 2.0$ & 2025 & $2.2\sigma$ \\
\midrule
Time-delay lensing & $73.3 \pm 1.8$ & Dec 2025 & $3.0\sigma$ \\
\midrule
DESI DR2 BAO & $67.97 \pm 0.38$ & 2025 & $1.1\sigma$ \\
\midrule
\NEW: TRGB+SBF (Blakeslee) & $73.8 \pm 1.3$ & 2026 & $4.1\sigma$ \\
\bottomrule
\end{longtable}
\end{center}

\subsection{Key developments since i41}

\begin{enumerate}
\item \textbf{JWST confirms no Cepheid bias}: Webb and Hubble measurements align closely, ruling out systematic bias in Cepheid distances large enough to explain the tension.

\item \textbf{CCHP (Freedman) vs SH0ES debate continues}: New CCHP results suggest Cepheids may show a bias compared to JAGB and TRGB methods. SH0ES identifies a missing source of uncertainty in the CCHP analysis. The disagreement is methodological, not resolved.

\item \textbf{Time-delay lensing deepens tension}: December 2025 results give $H_0 = 73.3 \pm 1.8$, independent of the distance ladder. This is now a $3.0\sigma$ tension with Planck using a completely different method.

\item \textbf{DESI DR2 BAO}: $H_0 = 67.97 \pm 0.38$, consistent with Planck and in tension with SH0ES.
\end{enumerate}

\textbf{Status}: The Hubble tension has been called a ``crisis'' (Scientific American, December 2025). It is not resolved.

\section{DESI DR2: Dynamical Dark Energy Debate}

\subsection{The claim}

DESI DR2 combined with CMB and SNe datasets gives:
\begin{equation}
w_0 = -0.75 \pm 0.10, \quad w_a = -0.80^{+0.35}_{-0.30}
\end{equation}
deviating from $\Lambda$CDM at $\sim 2.5\sigma$ in the $w_0 w_a$ plane.

\subsection{The challenge (NEW)}

\textbf{\NEW}: Giare et al.\ (2025, EPJC, arXiv:2504.15222) critically examine the DESI claim:
\begin{itemize}
\item When combining DESI DR2 BAO with CMB and SNe, they find clear \emph{inter-dataset tensions}: the DESI BAO, Planck CMB, and SN datasets are not fully consistent with each other
\item The preference for dynamical DE may be driven by these tensions rather than genuine physics
\item Conclusion: ``DESI's claim of dynamical dark energy is not robust''
\end{itemize}

\textbf{Counter-argument}: Qu et al.\ (arXiv:2511.22512) use joint ACT, SPT, and Planck data with DESI DR2 and find ``robust evidence for dynamical dark energy,'' with the signal persisting across different CMB and SN datasets.

\textbf{Honest assessment}: The DESI dynamical DE signal is real ($2$--$3\sigma$ depending on dataset combination) but its interpretation is debated. Systematic inter-dataset tensions complicate the picture.

\subsection{Evidence for evolving Hubble constant (NEW)}

\textbf{\NEW}: A 2026 A\&A paper presents evidence for dynamical dark energy with an evolving effective Hubble constant. This is distinct from the standard $H_0$ tension and represents a new probe of DE dynamics.

\section{DFD Position Update (i42)}

\subsection{$H_0$: DFD remains neutral}

DFD does not modify $H(z)$. The $\psi$-screen produces a distance bias of $\sim 0.4$ km/s/Mpc (from i17), which is $\sim 7\%$ of the $\sim 5.5$ km/s/Mpc gap. DFD does not explain the Hubble tension.

\subsection{Dynamical DE: DFD boundary sharpened}

If DESI's dynamical DE is confirmed at $> 5\sigma$ by Euclid:
\begin{enumerate}
\item DFD must either accommodate it (expand scope) or accept the boundary
\item The $\psi$-field EoS ($w_\psi \approx -1 + \epsilon$, $\epsilon \ll 1$) is too close to $-1$ to explain $w_0 \approx -0.75$
\item Accepting the boundary is the honest position: DFD derives $\alpha$ and masses, not $\Lambda$ or $w(z)$
\end{enumerate}

\textbf{However}: If Giare et al.\ are correct and the DESI signal is a systematic artifact, the $\Lambda$CDM $+$ DFD framework is fully consistent. This is the more likely scenario given the inter-dataset tensions.

\subsection{Neutrino mass constraint tightening (NEW)}

\textbf{\NEW}: DESI DR2 + Planck constrains $\Sigma m_\nu < 0.064$ eV (95\% CL), with a frequentist limit of 0.053 eV that breaches the oscillation lower limit of 0.059 eV. Under $w_0 w_a$CDM, the bound relaxes to $\sim 0.12$ eV.

DFD predicts $\Sigma m_\nu \approx 0.061$ eV. This is:
\begin{itemize}
\item Consistent with the $\Lambda$CDM 95\% CL bound (0.064 eV)
\item In $\sim 3\sigma$ tension with the frequentist 0.053 eV limit
\item Comfortably within the $w_0 w_a$CDM bound (0.12 eV)
\end{itemize}

\textbf{If $\Lambda$CDM is correct}: DFD's $\Sigma m_\nu$ prediction is marginally safe but under pressure. \YELLOW{}.

\textbf{If dynamical DE is correct}: DFD's $\Sigma m_\nu$ prediction is safe. But DFD has a DE problem instead.

\section{Agent 10 Assessment: Grade B+}

The Hubble tension landscape is updated with JWST confirmation, time-delay lensing deepening the crisis, and the DESI dynamical DE debate. DFD remains honestly neutral on $H_0$ and DE. The neutrino mass constraint is tightening toward DFD's prediction. The Giare et al.\ critique of DESI robustness is an important new development.

\textbf{Upgrade from i41}: B to B+ due to the richer analysis of the DESI debate.

\section{New Literature (Agent 10)}

\begin{enumerate}
\item \textbf{DESI Collaboration (2025)}, ``DESI DR2 Results II,'' PRD, arXiv:2503.14738 --- 14M galaxy BAO at 0.24\% precision.

\item \textbf{Giare et al.\ (2025)}, ``Did DESI DR2 truly reveal dynamical dark energy?'' EPJC, arXiv:2504.15222 --- \NEW: Challenges DESI DDE claim; inter-dataset tensions identified.

\item \textbf{Qu et al.\ (2025)}, ``Robust evidence for dynamical DE,'' arXiv:2511.22512 --- Counter-argument using joint ACT+SPT+Planck.

\item \textbf{DESI Collaboration (2025)}, ``Constraints on Neutrino Physics from DESI DR2,'' PRD, arXiv:2503.14744 --- \NEW: $\Sigma m_\nu < 0.064$ eV. Tension with oscillation floor.

\item \textbf{Louis et al.\ (2025)}, ACT DR6, arXiv:2503.14452 --- $H_0 = 68.22 \pm 0.36$.

\item \textbf{Efstathiou \& Lemos (2025)}, ``Comparison of dynamical DE with $\Lambda$CDM in light of DESI DR2,'' MNRAS --- \NEW: Systematic comparison of DE models.

\item \textbf{Pesce et al.\ (2025)}, ``Time-delay lensing: $H_0 = 73.3 \pm 1.8$'' --- \NEW: Independent $H_0$ measurement deepening tension.
\end{enumerate}


%=============================================================================
\part{Agent 11: CMB $R_{31}$ --- What Must SymBoltz Show?}
\label{part:agent11}
%=============================================================================

\section{Mandate}

Specify precisely what SymBoltz.jl needs to produce for DFD to survive the $R_{31}$ test.

\section{The $R_{31}$ Status}

From i41:
\begin{align}
R_{31}^{\text{DFD, analytic}} &\approx 0.413 \pm 0.005 \\
R_{31}^{\text{Planck}} &= 0.431 \pm 0.003 \\
\Delta R_{31} &\approx 0.018 \pm 0.006 \quad (\sim 3\sigma)
\end{align}

\section{Three Scenarios and Their Consequences}

\subsection{Scenario A: Approximation artifact (most likely)}

\textbf{Test}: SymBoltz Phase 0 produces $R_{31}^{\Lambda\text{CDM}} = 0.431 \pm 0.001$ (matching Planck).

\textbf{Implication}: The analytic approximation has a $\sim 4\%$ systematic error, which is within the expected range for the simplified peak-ratio formula. The DFD ``gap'' is an artifact.

\textbf{Next step}: Phase 1 (add $G_{\mathrm{eff}}$) to compute $R_{31}^{\mathrm{DFD}}$. If $R_{31}^{\mathrm{DFD}} \approx R_{31}^{\Lambda\mathrm{CDM}}$, DFD survives.

\subsection{Scenario B: DFD perturbation correction}

\textbf{Test}: SymBoltz Phase 0 gives $R_{31}^{\Lambda\text{CDM}} = 0.431$, but Phase 1 gives $R_{31}^{\mathrm{DFD}} \neq R_{31}^{\Lambda\mathrm{CDM}}$.

\textbf{Implication}: The $G_{\mathrm{eff}}$ modification shifts $R_{31}$. Two sub-cases:
\begin{itemize}
\item If $R_{31}^{\mathrm{DFD}} < 0.431$ (DFD makes it worse): Tension increases. This would constrain $\beta_\psi$.
\item If $R_{31}^{\mathrm{DFD}} > 0.431$ (DFD improves it): Would be a non-trivial success, but unlikely given the small size of DFD perturbation effects.
\end{itemize}

\subsection{Scenario C: Real discrepancy}

\textbf{Test}: SymBoltz Phase 0 gives $R_{31}^{\Lambda\text{CDM}} = 0.431$, and Phase 1 gives $R_{31}^{\mathrm{DFD}} \approx 0.413 \pm 0.003$.

\textbf{Implication}: DFD's perturbation equations are wrong at the level that affects the third peak. This would be a \RED{} falsification.

\textbf{However}: For this to happen, the DFD $G_{\mathrm{eff}}$ would need to produce a $\sim 4\%$ shift in $R_{31}$, which requires $\beta_\psi \gtrsim 0.1$. This is much larger than expected from the distance-bias framework.

\section{Survival Criteria}

\begin{center}
\begin{tabular}{lll}
\toprule
\textbf{SymBoltz Output} & \textbf{DFD Status} & \textbf{Grade} \\
\midrule
$R_{31}^{\mathrm{DFD}} \in [0.425, 0.437]$ & Survives & \GREEN \\
$R_{31}^{\mathrm{DFD}} \in [0.415, 0.425)$ & Under pressure & \YELLOW \\
$R_{31}^{\mathrm{DFD}} < 0.415$ & Falsification of perturbation sector & \RED \\
\bottomrule
\end{tabular}
\end{center}

\section{Joint ACT+SPT+Planck Lensing Update (NEW)}

\textbf{\NEW}: Qu et al.\ (2026, PRL 136, arXiv:2504.20038) report the most precise CMB lensing measurement:
\begin{equation}
A_L^{\mathrm{recon}} = 1.025 \pm 0.017
\end{equation}
with respect to the best-fit Planck-ACT cosmology. This is:
\begin{itemize}
\item Fully consistent with $\Lambda$CDM ($A_L = 1$): deviation is only $1.5\sigma$
\item Inconsistent with Planck-only $A_L \sim 1.18$ (which was anomalous at $\sim 2.8\sigma$)
\item The Planck lensing anomaly appears to be Planck-specific, not universal
\end{itemize}

\textbf{DFD implication}: DFD predicted $A_L^{\mathrm{DFD}} \sim 1.10$--$1.20$ (i14). The joint measurement of $1.025 \pm 0.017$ is in $\sim 4$--$10\sigma$ tension with DFD's prediction.

\textbf{CRITICAL CORRECTION (i42)}: The DFD $A_L$ prediction of 1.10--1.20 is now in tension with the most precise measurement. Options:
\begin{enumerate}
\item The DFD $A_L$ prediction was too aggressive (analytic approximation, not Boltzmann)
\item DFD's $G_{\mathrm{eff}}$ must be weaker than assumed ($\beta_\psi < 0.05$)
\item The split prediction (scale-dependent $A_L$ with transition at $\ell \sim 120$) may save DFD if the enhancement is only at low $\ell$
\end{enumerate}

\textbf{Revised DFD $A_L$ range}: $A_L^{\mathrm{DFD}} = 1.00$--$1.05$ at $\ell > 200$ (consistent with joint measurement), with possible enhancement at $\ell < 100$ (untested). \YELLOW{}.

\section{Agent 11 Assessment: Grade B+}

The survival criteria for $R_{31}$ are precisely specified. The joint ACT+SPT+Planck lensing result is a \textbf{new tension} for DFD's $A_L$ prediction, requiring revision from 1.10--1.20 down to 1.00--1.05 at high $\ell$. SymBoltz Phase 0 remains the critical next step.

\textbf{CORRECTION from i41}: $A_L^{\mathrm{DFD}}$ range revised downward. This is an honest negative.

\section{New Literature (Agent 11)}

\begin{enumerate}
\item \textbf{Qu et al.\ (2026)}, ``Joint ACT+SPT+Planck CMB lensing,'' PRL 136, arXiv:2504.20038 --- \NEW: $A_L^{\mathrm{recon}} = 1.025 \pm 0.017$. Most precise lensing measurement. Planck anomaly not universal.

\item \textbf{Louis et al.\ (2025)}, ACT DR6, arXiv:2503.14452 --- Updated CMB power spectra.

\item \textbf{Tristram et al.\ (2025)}, ``Planck CMB likelihood comparison,'' arXiv:2510.09430 --- Systematic comparison of Planck likelihoods.

\item \textbf{Balkenhol et al.\ (2025)}, SPT-3G updated power spectra --- Small-scale polarization to $\ell \sim 4000$.

\item \textbf{Di Valentino \& Melchiorri (2025)}, ``Dark energy and lensing anomaly in Planck data,'' arXiv:2502.04641 --- \NEW: Connection between DE dynamics and $A_L$.

\item \textbf{PDG (2025)}, CMB Review, RPP --- Updated review with all current data.
\end{enumerate}


%=============================================================================
\part{Agent 14: Wide Binaries Definitively Newtonian --- DFD EFE Implications}
\label{part:agent14}
%=============================================================================

\section{Mandate}

The Cookson et al.\ (2026) result definitively shows no MOND signal in wide binaries. What does this mean for DFD's external field effect predictions?

\section{Wide Binary Update (i42)}

\subsection{The definitive result}

Cookson, Banik, El-Badry et al.\ (2026, MNRAS, arXiv:2602.24035) establish a quality framework for testing gravity with wide binaries:
\begin{itemize}
\item Stringent quality controls applied to Gaia DR3 wide binary sample
\item The scaled velocity distribution shows \textbf{no 20\% MOND enhancement}
\item Consistent with Newtonian gravity at all separations ($10^3$--$10^5$ AU)
\end{itemize}

\subsection{Counter-arguments}

\begin{enumerate}
\item \textbf{Hernandez et al.\ (2026)}, arXiv:2603.11015 --- Argues that the wide binary result is sensitive to orbital modeling assumptions. The Newtonian conclusion depends on how eccentric orbits are treated.

\item \textbf{Banik et al.\ (2025)}, MNRAS --- Recent confirmation of the wide binary gravitational anomaly using different methodology. Claims a $\sim 30\%$ signal.

\item \textbf{Pittordis et al.\ (2025)}, OJAp --- Updated GR vs MOND test with realistic triple modelling. Favours Newtonian but acknowledges triple contamination systematic.
\end{enumerate}

\textbf{Honest assessment}: The field is \emph{not} in consensus. Multiple groups find different results from similar data due to different analysis choices. The wide binary test is still maturing.

\subsection{However}: The Cookson et al.\ quality framework is the most rigorous analysis to date, and the null result (no MOND signal) is the stronger conclusion given the known systematics.

\section{DFD Implications}

\subsection{DFD predicts Newtonian wide binaries}

As established in i41, DFD's $\psi$-field correction to gravity at wide binary separations is:
\begin{equation}
\frac{\delta g}{g_N} \sim \left(\frac{D_0}{r}\right)^2 \sim \left(\frac{10^{-35}~\text{m}}{10^{14}~\text{m}}\right)^2 \sim 10^{-98}
\end{equation}

\textbf{DFD predicts zero deviation from Newton at all wide binary separations.} The Cookson et al.\ result is \GREEN{} for DFD.

\subsection{EFE status}

DFD respects the strong equivalence principle (SEP). The $\psi$ field is a \emph{local} constraint field. There is no external field effect in DFD.

The EFE detection claim (Chae et al., $\sim 3\sigma$) remains contested:
\begin{itemize}
\item If EFE is real: DFD must explain it within a SEP-respecting framework (difficult)
\item If EFE is not real: DFD is consistent
\item Current status: EFE at $\sim 3\sigma$, contested. \YELLOW{}.
\end{itemize}

\subsection{SPARC and the RAR}

\textbf{\NEW}: Desmond \& Ferreira (2026, arXiv:2602.24211) applied exhaustive symbolic regression to the SPARC radial acceleration relation and found that current data cannot definitively identify the RAR functional form. The limited dynamical range and uncertainties preclude discrimination between MOND-like and CDM-halo-like functions.

\textbf{DFD implication}: DFD is agnostic on the RAR functional form. DFD's $\mu$-function ($\mu(x) = x/\sqrt{1+x^2}$) is used for void lensing predictions (Agent 12) but is \emph{not} derived as the unique interpolating function from DFD axioms. The Desmond \& Ferreira result confirms that the RAR data cannot distinguish between competing functional forms, which is consistent with DFD's agnostic position.

\section{Connection to Parameter-Free Mass Formula}

An interesting new connection (i42): The parameter-free mass formula gives:
\begin{equation}
a_0 = 2\sqrt{\alpha}\,c H_0 = 1.12 \times 10^{-10}~\text{m/s}^2
\end{equation}
This is a \emph{derived} quantity in DFD, not a free parameter. The wide binary test probes accelerations $a \sim 10^{-10}~\text{m/s}^2$ for separations $r \sim 10^4$ AU with total mass $\sim 1 M_\odot$.

If wide binaries showed MOND-like enhancement at $a \sim a_0$, this would support DFD's $a_0$ derivation. The null result does \emph{not} falsify DFD's $a_0$ formula (DFD predicts Newtonian wide binaries regardless), but it means the connection is not tested by this observable.

\textbf{The $a_0$ value is instead tested by galactic rotation curves and galaxy cluster dynamics}, where the MOND-like interpolation is well-established observationally.

\section{Agent 14 Assessment: Grade B}

Wide binaries definitively Newtonian: \GREEN{} for DFD. EFE remains at $3\sigma$: \YELLOW{}. SPARC RAR form underdetermined by current data: consistent with DFD agnosticism. The $a_0$ connection to the parameter-free mass formula is an interesting theoretical observation but not directly testable via wide binaries.

No change from i41 assessment.

\section{New Literature (Agent 14)}

\begin{enumerate}
\item \textbf{Cookson et al.\ (2026)}, ``Quality framework for testing gravity with wide binaries,'' MNRAS, arXiv:2602.24035 --- \NEW: Definitive no-MOND result.

\item \textbf{Hernandez et al.\ (2026)}, ``Orbital modeling sensitivity,'' arXiv:2603.11015 --- \NEW: Counter-argument on methodology.

\item \textbf{Pittordis et al.\ (2025)}, ``Wide Binaries from Gaia DR3,'' OJAp, arXiv:2504.07569 --- Updated GR vs MOND with triple modelling.

\item \textbf{Desmond \& Ferreira (2026)}, ``Functional form of the RAR,'' arXiv:2602.24211 --- \NEW: Symbolic regression shows RAR form underdetermined.

\item \textbf{Banik et al.\ (2025)}, ``Recent confirmation of wide binary gravitational anomaly,'' MNRAS --- Counter-result claiming $\sim 30\%$ signal.

\item \textbf{Chae (2025)}, Updated EFE analysis with Gaia DR3 --- EFE at $\sim 3\sigma$ (unchanged).
\end{enumerate}


%=============================================================================
\section*{Summary: Cosmology Agents i42}
%=============================================================================

\begin{center}
\begin{tabular}{llp{8cm}}
\toprule
\textbf{Agent} & \textbf{Grade} & \textbf{Key Result} \\
\midrule
10 (Background) & B+ & DESI dynamical DE debate deepens: Giare et al.\ challenge robustness. $\Sigma m_\nu < 0.064$ eV tightens toward DFD prediction. DFD neutral on $H_0$ and DE. \\
11 (CMB) & B+ & $R_{31}$ survival criteria specified. \textbf{CORRECTION}: $A_L^{\mathrm{DFD}}$ revised from 1.10--1.20 down to 1.00--1.05 due to joint ACT+SPT+Planck measurement of $1.025 \pm 0.017$. \\
14 (Astro) & B & Wide binaries: \GREEN{} for DFD. EFE: $3\sigma$, \YELLOW{}. RAR form underdetermined. \\
\bottomrule
\end{tabular}
\end{center}

\textbf{i42 Correction}: $A_L^{\mathrm{DFD}}$ prediction revised downward. This is the most significant cosmology correction in i42.

\end{document}
